% Author addition prepared after the retained Phase A results and before
% Phase B numerical execution.  No empirical outcome is asserted here.
\subsection{A bounded preliminary information budget}
\label{sec:budgetedtail119}
Resolving fewer innovation locations need not reduce complete decision
work.  Tail localization, several refused resolutions, and eventual repair
can outweigh the saved quadrature.  We therefore introduce a second
service whose preliminary work is bounded before the original complete
solver is invoked.  Phase B freezes this algorithm after the retained
Phase A outcomes; its acquisition, conventional controls, reference
intervals, and timing processes constitute a separate experiment.

Fix the original state, a feasible realized actor proposal, the original
risk level $\alpha\in\{1/2,3/4\}$, and the implemented terminal policy.
Two floating endpoint comparisons propose a boundary cut when possible.
Only the remaining rows receive six floating bisection steps.  These
signs have no certification authority: the resulting feasible cut is
rounded to the declared dyadic grid, and two outward original-continuation
values construct the defect $\kappa$ in
\eqref{eq:taildefect119}.  This defect is retained at every resolution.
The service omits exactly zero-length tails, attempts two and then four
equal cells per positive-length tail, and accepts only a certificate
meeting both original targets.  Otherwise it invokes the frozen
128-bin complete service from the same already acquired proposal,
including its original repair, refinement, and rational fallback.

\begin{proposition}[Original-policy qualification with bounded preliminary work]
\label{prop:budgetedtail119}
Under the outward-oracle and original-policy assumptions of
Theorem~\ref{thm:investmenttail119}, the stated service returns only
certificates valid for the original investment objective.  Its preliminary
work uses at most
\begin{equation}\label{eq:budgetedtailcount119}
 N_{\rm outward}\le 30,
 \qquad N_{\rm floating}\le16
\end{equation}
terminal point-value payloads per proposed action before acceptance or
entry to the frozen complete service.  An outward payload includes the
original value enclosure and the valid supporting information required
by the action certificate.  Calls repeated at the same location are
counted again.

More precisely, if the selected density has $p\in\{1,2\}$ nonempty
tails, the outward bound is $2+14p$; an acceptance at the first
resolution uses at most $2+5p$.  A floating endpoint exit uses only four
floating values.  The two outward defect values are acquired once and
shared across the preliminary resolutions.  If the underlying oracle
calls and frozen complete service terminate, adding this finite prelude
preserves that termination and the original loss and width targets.
\end{proposition}

These are bounds on added terminal information, not on elapsed time or
on the complete internal work of a terminal solve.  Actor acquisition,
the payloads' numerical construction, rejected attempts, density and
certificate algebra, and the complete invoked fallback remain paid work.
Early acceptance can avoid a conventional repair and can therefore
change the returned action.  Its original economic cost must be compared
using its own reference interval; no zero-cost or improvement conclusion
follows from the point-count bound.

The masked-tail ablation applies the same exact zero-mass omission and
floating proposal rule while retaining all six resolutions
$2,4,8,16,32,64$.  It separates the effect of removing zero-mass work
from the additional decision to cap preliminary effort.  Every
representation, including the conventional hinge encoding of ReLU,
receives the same service and information access.
